\section{Related Work}

\subsection{Weight-Space Learning}

\textbf{Hyper-Representations and Meta-Learning.} Sch\"urholt et al.~\cite{schurholt2022hyper} introduce hyper-representations as learned embeddings of neural network weights for transfer learning tasks, demonstrating that weight-based features generalize across model zoos. Our work complements this by validating the layer-wise tokenization assumption underlying such embeddings.

\textbf{Model Stitching and Layer Analysis.} Lenc \& Vedaldi~\cite{lenc2015understanding} show that layers can be analyzed semi-independently while preserving model behavior, supporting the hypothesis that cross-layer dependencies are not always critical. We extend this by empirically measuring signal preservation under layer-wise processing on property prediction tasks.

\textbf{Weight Statistics and Property Prediction.} Unterthiner et al.~\cite{unterthiner2020predicting} predict model accuracy from weight statistics (norms, spectral properties), achieving moderate correlation. Building on this validation of weight-space signal, we compare statistical aggregation directly against transformer-based tokenization, revealing that on small datasets (100 models), both approaches achieve similar performance (80\% $\pm$ 4.2\% accuracy).

\textbf{Neural Functionals.} Recent work processes weights as graph structures or functional mappings \cite{navon2023equivariant,zhou2022neural}. While promising, these methods introduce architectural complexity. Our contribution isolates tokenization quality independent of backbone architecture choice.

\subsection{Sequence Modeling on Structured Data}

\textbf{Transformers for Point Clouds and Graphs.} Point cloud transformers \cite{zhao2021point} and graph transformers \cite{dwivedi2021generalization} process variable-length structured data via attention mechanisms. Weight-space learning shares similar challenges—variable layer dimensions, permutation symmetries—but operates on learned parameter distributions rather than geometric coordinates.

\textbf{Positional Encodings for Hierarchical Data.} Hierarchical position encodings \cite{shaw2018self} capture depth information in trees and nested structures. We adapt this for weight tokens where layer depth indicates network position, enabling transformers to attend across architectural boundaries.

\subsection{Meta-Learning on Model Zoos}

\textbf{Architecture Search and Transfer Learning.} Neural Architecture Search (NAS) methods \cite{zoph2017neural} predict model performance from architectural descriptors. Weight-space learning extends this by using \textit{learned parameters} as features rather than discrete architecture choices, enabling fine-grained property prediction post-training.

\textbf{Model Merging and Task Arithmetic.} Recent work \cite{wortsman2022model,ilharco2023editing} merges pretrained models via weight averaging or interpolation. Our tokenization validation is prerequisite for learning \textit{how} to merge—embeddings must preserve model properties to predict merge outcomes.

\textbf{Backdoor Detection in Model Weights.} Spectral signatures \cite{wang2019neural} and statistical anomalies \cite{tran2018spectral} detect backdoored models from weight distributions. Our work validates that layer-wise tokenization retains sufficient signal for such anomaly detection tasks, with 80\% family classification suggesting discriminative capacity for finer-grained backdoor patterns.

\subsection{Positioning}

Existing weight-space learning methods demonstrate task-specific success but lack systematic evaluation of tokenization strategies. We provide a controlled experiment comparing layer-wise processing approaches, showing that:

\begin{enumerate}
\item \textbf{Tokenization preserves signal}: 80\% $\pm$ 4.2\% accuracy significantly exceeds random baseline (25\%) and gate threshold (60\%)
\item \textbf{Simple baselines matter}: Per-layer statistics match transformer performance on small datasets, establishing rigorous comparison point
\item \textbf{Family separability dominates}: On 100-model zoo, architecture-specific patterns (BatchNorm, attention weights, kernel structures) are detectable via local layer statistics without requiring cross-layer reasoning
\end{enumerate}

This validation enables future work on larger datasets, harder tasks, and more sophisticated weight-processing backbones with confidence that the foundational layer-wise assumption holds.
